\chapter*{Bảng thuật ngữ}
\addcontentsline{toc}{chapter}{Bảng thuật ngữ}
\label{glossary}

Báo cáo sử dụng một số thuật ngữ chuyên ngành thuộc lĩnh vực blockchain, tài chính số và kiến trúc phần mềm. Phần lớn các thuật ngữ này đã được dịch hoặc diễn giải bằng tiếng Việt ngay trong nội dung từng chương; bảng dưới đây tổng hợp lại những thuật ngữ thông dụng nhất để tiện tra cứu trong suốt quá trình đọc báo cáo.

\renewcommand{\arraystretch}{1.3}
\small
\begin{longtable}{|L{3.8cm}|L{11.2cm}|}
\hline
\textbf{Thuật ngữ} & \textbf{Giải thích} \\
\hline
\endfirsthead

\multicolumn{2}{c}{\textit{(Bảng thuật ngữ - tiếp theo)}} \\
\hline
\textbf{Thuật ngữ} & \textbf{Giải thích} \\
\hline
\endhead

\hline
\multicolumn{2}{r}{\textit{(còn tiếp ở trang sau)}} \\
\endfoot

\hline
\endlastfoot

Adapter & Lớp trung gian trong backend, chuyển dữ liệu từ định dạng riêng của một nhà cung cấp bên ngoài sang định dạng nội bộ thống nhất của hệ thống. \\
\hline

AI Assistant (Trợ lý AI) & Thành phần sử dụng mô hình ngôn ngữ lớn để tóm tắt và diễn giải dữ liệu ví, token hoặc tín hiệu wash trading dựa trên ngữ cảnh đã được hệ thống chuẩn hóa. \\
\hline

API (Application Programming Interface) & Giao diện lập trình ứng dụng, tập hợp quy tắc để các thành phần phần mềm giao tiếp và trao đổi dữ liệu với nhau. \\
\hline

ATH / ATL (All-Time High / All-Time Low) & Mức giá cao nhất và thấp nhất mà một token từng đạt được tính đến thời điểm hiện tại. \\
\hline

Backend & Phần máy chủ của ứng dụng, đảm nhiệm xác thực, xử lý nghiệp vụ, truy cập cơ sở dữ liệu và tích hợp với các nhà cung cấp dữ liệu bên ngoài. \\
\hline

Blockchain & Chuỗi khối, một dạng sổ cái phân tán lưu trữ dữ liệu giao dịch một cách công khai, có thể kiểm chứng và khó bị thay đổi. \\
\hline

Cache / Caching (Lưu trữ đệm) & Cơ chế lưu lại dữ liệu đã truy xuất để tái sử dụng cho các lần yêu cầu sau, giúp giảm số lượt gọi đến nhà cung cấp bên ngoài và cải thiện tốc độ phản hồi. \\
\hline

DeFi (Decentralized Finance) & Tài chính phi tập trung, các dịch vụ tài chính như cho vay, giao dịch hay tạo thanh khoản được vận hành trực tiếp trên blockchain, không qua trung gian truyền thống. \\
\hline

DEX (Decentralized Exchange) & Sàn giao dịch phi tập trung, nơi người dùng hoán đổi token trực tiếp thông qua hợp đồng thông minh thay vì thông qua một đơn vị trung gian. \\
\hline

Drawdown & Mức suy giảm giá trị tài sản của một ví so với đỉnh giá trị gần nhất, dùng để đánh giá mức độ rủi ro. \\
\hline

DTO (Data Transfer Object) & Đối tượng dữ liệu được định nghĩa để truyền dữ liệu giữa các lớp của hệ thống theo một cấu trúc thống nhất. \\
\hline

Endpoint & Một địa chỉ cụ thể trong API mà frontend hoặc dịch vụ khác gửi yêu cầu đến để lấy hoặc gửi dữ liệu. \\
\hline

ERD (Entity Relationship Diagram) & Sơ đồ thực thể mối quan hệ, mô tả các nhóm bảng chính và mối liên hệ giữa chúng trong cơ sở dữ liệu. \\
\hline

FDV (Fully Diluted Valuation) & Định giá pha loãng hoàn toàn, giá trị vốn hóa ước tính của một token nếu toàn bộ nguồn cung tối đa được lưu hành. \\
\hline

Frontend & Phần giao diện người dùng của ứng dụng, đảm nhiệm hiển thị dữ liệu, tiếp nhận thao tác và trực quan hóa thông tin. \\
\hline

Full-stack & Kiến trúc ứng dụng bao gồm cả phần giao diện (frontend) và phần máy chủ (backend) được phát triển trong cùng một hệ thống. \\
\hline

GCN / GAT / GraphSAGE & Các kiến trúc mạng nơ-ron đồ thị (Graph Neural Network) dùng để học đặc trưng và phát hiện quan hệ bất thường trên dữ liệu dạng đồ thị ví--giao dịch. \\
\hline

Holder & Địa chỉ ví đang nắm giữ một loại token nhất định. \\
\hline

Lazy Loading (Tải dữ liệu động) & Kỹ thuật chỉ tải dữ liệu khi thực sự có nhu cầu sử dụng, thay vì tải toàn bộ dữ liệu ngay từ đầu. \\
\hline

LLM (Large Language Model) & Mô hình ngôn ngữ lớn, được sử dụng để diễn giải các chỉ số đã tính toán thành nội dung bằng ngôn ngữ tự nhiên. \\
\hline

Monolith & Kiến trúc backend được triển khai như một ứng dụng thống nhất duy nhất, nhưng bên trong được chia thành nhiều lớp và miền nghiệp vụ riêng biệt. \\
\hline

Monorepo & Cách tổ chức mã nguồn trong đó nhiều thành phần của dự án như frontend, backend và tài liệu kỹ thuật cùng nằm trong một kho mã nguồn duy nhất. \\
\hline

NFT (Non-Fungible Token) & Token không thể thay thế, đại diện cho một tài sản số có tính duy nhất trên blockchain. \\
\hline

On-chain / Off-chain & On-chain là dữ liệu được ghi nhận trực tiếp trên blockchain; off-chain là dữ liệu nằm ngoài blockchain, không được ghi nhận trực tiếp trên chuỗi. \\
\hline

ORM (Object-Relational Mapping) & Công cụ ánh xạ giữa cấu trúc bảng trong cơ sở dữ liệu quan hệ với các đối tượng trong mã nguồn, giúp truy vấn dữ liệu an toàn hơn về kiểu dữ liệu. \\
\hline

Partial-range cache reuse & Kỹ thuật tái sử dụng phần dữ liệu đã đồng bộ trong một khoảng thời gian, chỉ truy xuất bổ sung phần dữ liệu còn thiếu khi người dùng thay đổi khoảng thời gian phân tích. \\
\hline

PnL (Profit and Loss) & Lãi/lỗ, chỉ số phản ánh kết quả tăng hoặc giảm giá trị tài sản của một ví theo thời gian. \\
\hline

Pool (Pool thanh khoản) & Quỹ thanh khoản chứa cặp tài sản, dùng để thực hiện giao dịch hoán đổi trên các sàn phi tập trung. \\
\hline

Provider (Nhà cung cấp dữ liệu) & Nhà cung cấp dữ liệu bên thứ ba, ví dụ Helius, CoinGecko hoặc Birdeye, nơi hệ thống truy xuất dữ liệu on-chain hoặc dữ liệu thị trường. \\
\hline

Rate limit (Giới hạn truy vấn) & Giới hạn về số lượng hoặc tần suất yêu cầu mà một API cho phép trong một khoảng thời gian nhất định. \\
\hline

Request coalescing & Cơ chế gộp nhiều yêu cầu dữ liệu giống nhau phát sinh đồng thời thành một lần gọi duy nhất đến nhà cung cấp. \\
\hline

RPC (Remote Procedure Call) & Giao thức gọi hàm từ xa, được dùng để tương tác trực tiếp với node của mạng lưới Solana. \\
\hline

Schema (Lược đồ dữ liệu) & Mô tả cấu trúc các bảng, cột và mối quan hệ trong cơ sở dữ liệu. \\
\hline

Sequence Diagram (Sơ đồ tuần tự) & Sơ đồ mô tả trình tự tương tác giữa các thành phần của hệ thống theo thời gian. \\
\hline

Shard / Sharding & Kỹ thuật phân chia dữ liệu hoặc cấu hình, ví dụ danh sách địa chỉ theo dõi, thành nhiều nhóm nhỏ hơn để dễ quản lý và tránh vượt giới hạn của nhà cung cấp. \\
\hline

Sparkline & Biểu đồ nhỏ, đơn giản, dùng để thể hiện nhanh xu hướng biến động của một chỉ số theo thời gian. \\
\hline

Stablecoin & Loại token có giá trị được neo theo một tài sản ổn định, phổ biến nhất là đồng USD. \\
\hline

Swap & Giao dịch hoán đổi trực tiếp giữa hai loại token thông qua giao thức phi tập trung. \\
\hline

Token & Đơn vị tài sản số được phát hành và giao dịch trên blockchain, có thể đại diện cho tiền mã hóa, quyền sở hữu hoặc một loại tài sản khác. \\
\hline

TTL (Time To Live) & Khoảng thời gian một dữ liệu cache được xem là còn hợp lệ trước khi hệ thống cần cập nhật lại từ nhà cung cấp. \\
\hline

Use Case (Tình huống sử dụng) & Mô tả một tương tác cụ thể giữa tác nhân, có thể là người dùng hoặc hệ thống bên ngoài, với hệ thống đang xây dựng. \\
\hline

Wallet (Ví) & Địa chỉ dùng để lưu trữ, quản lý và thực hiện giao dịch tài sản số trên blockchain. \\
\hline

Wash trading & Hành vi giao dịch được phối hợp thực hiện nhằm tạo ra khối lượng hoặc mức độ hoạt động giả cho một tài sản. \\
\hline

Watchlist (Danh sách theo dõi) & Danh sách token hoặc ví được người dùng lưu lại để tiện theo dõi. \\
\hline

Web3 & Thế hệ ứng dụng web được xây dựng trên nền tảng blockchain, hướng đến việc phi tập trung hóa dữ liệu và quyền kiểm soát. \\
\hline

Webhook & Cơ chế cho phép một dịch vụ tự động gửi dữ liệu đến một địa chỉ đã đăng ký ngay khi có sự kiện xảy ra, thay vì phải chủ động truy vấn liên tục. \\
\hline

\end{longtable}
\normalsize